\chapter{Phân tích yêu cầu và thiết kế hệ thống}
\section{Câu hỏi nghiên cứu}
Báo cáo được tổ chức quanh ba câu hỏi nghiên cứu chính:
\begin{enumerate}
    \item \textbf{RQ1:} \TFIDF{} kết hợp Linear SVM có đủ mạnh cho phân loại sentiment review TMĐT tiếng Việt trong phạm vi dữ liệu hiện có không?
    \item \textbf{RQ2:} Retrieval/RAG có thể tạo insight có bằng chứng từ kho review thay vì chỉ đưa ra nhãn cảm xúc đơn lẻ không?
    \item \textbf{RQ3:} Hybrid retrieval kết hợp BM25, Dense FAISS, RRF/rerank và aspect-aware rules có cải thiện khả năng bắt keyword/aspect complaint trong benchmark nhỏ không?
\end{enumerate}
\section{Yêu cầu bài toán}
Đề tài cần giải quyết bài toán phân tích review TMĐT tiếng Việt theo cách có thể kiểm chứng bằng dữ liệu và metric. Hệ thống phải nhận review để dự đoán sentiment, đồng thời nhận câu hỏi để truy xuất review liên quan và tạo câu trả lời có evidence. Người dùng cuối không cần biết chi tiết embedding hay FAISS; họ chỉ cần nhập review hoặc câu hỏi và xem kết quả được trình bày rõ ràng.

Yêu cầu quan trọng là không trộn input của sentiment với input của RAG. Một review sản phẩm là văn bản cần phân loại cảm xúc. Một câu hỏi RAG như “Khách hàng phàn nàn gì về giao hàng chậm?” không phải review, nên không nên đưa trực tiếp vào sentiment classifier như một review. Phiên bản giao diện cuối cùng tách hai chức năng này thành hai tab riêng.

\section{Yêu cầu chức năng và phi chức năng}
Yêu cầu chức năng gồm: nhập review và dự đoán sentiment; nhập câu hỏi RAG và truy xuất review liên quan; hiển thị answer có citation; lọc theo sentiment, category và rating; hiển thị cấu hình kỹ thuật và metric. Yêu cầu phi chức năng gồm: không cần API key, chạy local được, câu trả lời không bịa ngoài evidence, giao diện tách riêng sentiment và RAG, và artifact phải có đường dẫn rõ ràng để kiểm chứng.

\section{Input và output}
Input của Module 2 là review đã được làm sạch, rating và sentiment label. Output là model sentiment và các báo cáo metric. Input của Module 4 là corpus RAG, FAISS index, metadata và query người dùng. Output là danh sách review liên quan và câu trả lời có citation. Input của Module 6 là corpus/index hiện có và bộ query benchmark nhỏ. Output là kết quả hybrid retrieval và bảng so sánh keyword/aspect hit-rate.

\section{Kiến trúc tổng thể}
Hình \ref{fig:system-pipeline-tikz} mô tả pipeline theo luồng xử lý của hệ thống: dữ liệu thô đi qua Data Layer, sau đó đến sentiment model, analytics, retrieval/RAG, evaluation, Module 6 hybrid retrieval và cuối cùng là giao diện thử nghiệm. Cách vẽ tuần tự làm rõ output của module trước là cơ sở cho module sau.

\begin{figure}[H]
\centering
\begin{tikzpicture}[
node distance=0.65cm,
box/.style={rectangle, rounded corners, draw=teal!70!black, very thick, fill=teal!5, align=center, text width=0.82\textwidth, minimum height=0.9cm},
arrow/.style={-{Stealth[length=3mm]}, thick, draw=teal!70!black}
]
\node[box] (raw) {\textbf{Raw data}\\review, rating, product/category metadata};
\node[box, below=of raw] (m1) {\textbf{Module 1: Data Layer}\\clean text, deduplicate, map rating $\rightarrow$ sentiment, split, prepare RAG corpus};
\node[box, below=of m1] (m2) {\textbf{Module 2: Sentiment Classification}\\\TFIDF{} + Logistic Regression, \TFIDF{} + Linear SVM, select by validation \MacroFOne{}};
\node[box, below=of m2] (m3) {\textbf{Module 3: Review Analytics}\\sentiment distribution, \RatingSentiment{} crosstab, top negative reviews};
\node[box, below=of m3] (m4) {\textbf{Module 4: Retrieval/RAG}\\\MultilingualE{}, FAISS IndexFlatIP, metadata filtering, \EvidenceGrounded{} answer};
\node[box, below=of m4] (m5) {\textbf{Module 5: Evaluation \& Report}\\metrics summary, manual Precision@k, error analysis, report assembly};
\node[box, below=of m5] (m6) {\textbf{Module 6: Aspect-aware Hybrid Retrieval}\\BM25 + Dense FAISS + RRF/rerank + aspect-aware rules};
\node[box, below=of m6, fill=blue!4, draw=blue!60!black] (web) {\textbf{Giao diện thử nghiệm}\\FastAPI backend + frontend 3 tab: Sentiment, RAG, Kỹ thuật \& đánh giá};
\draw[arrow] (raw) -- (m1);
\draw[arrow] (m1) -- (m2);
\draw[arrow] (m2) -- (m3);
\draw[arrow] (m3) -- (m4);
\draw[arrow] (m4) -- (m5);
\draw[arrow] (m5) -- (m6);
\draw[arrow] (m6) -- (web);
\end{tikzpicture}
\caption{Kiến trúc pipeline NLP/RAG end-to-end của hệ thống.}
\label{fig:system-pipeline-tikz}
\end{figure}

\begin{table}[H]
\centering
\caption{Input/output chính của từng module}
\label{tab:module-io-contract}
\small
\begin{tabularx}{\textwidth}{L{0.18\textwidth}Y Y}
\toprule
Module & Input chính & Output chính \\
\midrule
Module 1 & Raw review, rating, metadata & Dữ liệu sentiment sau xử lý, corpus RAG và summary dữ liệu. \\
Module 2 & Train/valid/test sentiment & Model sentiment tốt nhất và bảng metric. \\
Module 3 & Dữ liệu sentiment/RAG đã xử lý & Phân phối sentiment, crosstab \RatingSentiment{}, top negative reviews và figures. \\
Module 4 & RAG corpus, FAISS index, metadata, query & Review liên quan, answer có citation và retrieval evaluation. \\
Module 5 & Metric Module 2/4 và error files & Final summary, discussion, conclusion và manual precision summary. \\
Module 6 & Corpus/index final và benchmark query & Hybrid retrieval output, benchmark reference và hướng nâng cấp. \\
Giao diện thử nghiệm & Model/index/artifact final & Giao diện 3 tab để chạy sentiment, RAG và xem kỹ thuật. \\
\bottomrule
\end{tabularx}
\end{table}

\clearpage
\section{Thiết kế module chi tiết}
Để báo cáo có cấu trúc gần với một bài báo kỹ thuật, mỗi module được mô tả theo bốn thành phần: mục tiêu, xử lý chính, artifact đầu ra và vai trò trong pipeline. Cách mô tả này giúp người đọc kiểm tra nhanh mối liên hệ giữa phương pháp, mã nguồn và kết quả thực nghiệm, đồng thời tránh trình bày hệ thống như một tập script rời rạc.

\begin{longtable}{L{0.13\textwidth}L{0.22\textwidth}L{0.34\textwidth}L{0.21\textwidth}}
\caption{Thiết kế chi tiết theo module của hệ thống}\label{tab:detailed-module-design}\\
\toprule
Module & Mục tiêu & Xử lý chính & Artifact/đầu ra \\
\midrule
Module 1 & Chuẩn hóa dữ liệu review. & Làm sạch text, chuẩn hóa rating, xử lý duplicate/conflict duplicate, ánh xạ rating sang sentiment và chia train/valid/test. & Dữ liệu sentiment sau xử lý, corpus RAG và summary dữ liệu. \\
Module 2 & Huấn luyện mô hình sentiment. & So sánh majority baseline, \TFIDF{} + Logistic Regression và \TFIDF{} + Linear SVM; chọn model bằng validation \MacroFOne{}. & Model tốt nhất, metric test, classification report và confusion matrix. \\
Module 3 & Phân tích dữ liệu review. & Tạo phân phối sentiment/rating, crosstab \RatingSentiment{}, top negative reviews và biểu đồ mô tả dữ liệu. & Bảng thống kê và figures phục vụ phân tích dữ liệu. \\
Module 4 & Truy xuất và tạo câu trả lời có bằng chứng. & Encode review bằng \MultilingualE{}, build FAISS IndexFlatIP, search top-k, filter/rerank nhẹ và generate answer theo template có citation. & FAISS index, metadata, index config, output retrieval/RAG. \\
Module 5 & Tổng hợp đánh giá. & Thu thập metric, manual retrieval evaluation, error analysis, discussion và conclusion. & Summary metric, Precision@k, phân tích lỗi và nội dung báo cáo. \\
Module 6 & Nâng cấp retrieval theo aspect. & Kết hợp BM25, dense FAISS, RRF fusion, optional reranker và rule-based aspect sentiment. & Benchmark hybrid retrieval và output demo rerank. \\
Giao diện & Kiểm thử pipeline end-to-end. & FastAPI load artifact, frontend gọi API, hiển thị sentiment, RAG answer, evidence cards và thông tin kỹ thuật. & Web demo 3 tab và API JSON. \\
\bottomrule
\end{longtable}

\section{Cấu trúc thư mục}
\begin{longtable}{p{0.28\textwidth}p{0.62\textwidth}}
\caption{Vai trò các thư mục chính trong mã nguồn thực nghiệm}\label{tab:project-folders}\\
\toprule
Thư mục/file & Vai trò \\
\midrule
\code{data/} & Chứa dữ liệu đã xử lý, corpus RAG, query evaluation và báo cáo dữ liệu. \\
\code{src/} & Chứa mã nguồn chính của Module 1--6. \\
\code{scripts/} & Chứa script chạy từng module, giao diện thử nghiệm và đóng gói. \\
\code{models/} & Chứa artifact mô hình sentiment và FAISS index/metadata. \\
\code{results/} & Chứa kết quả metric, dự đoán, output RAG, error analysis. \\
\code{figures/} & Chứa biểu đồ dùng trong báo cáo. \\
\code{report/} & Chứa skeleton và các đoạn báo cáo sinh từ Module 5. \\
\code{backend/} & FastAPI backend phục vụ sentiment và RAG API. \\
\code{frontend/} & Giao diện web gồm HTML, CSS, JavaScript. \\
\code{README_MODULE*.md} & Tài liệu giải thích từng module. \\
\code{requirements_*.txt} & Thư viện cần cài cho từng phần. \\
Manifest artifact & Danh sách artifact khi đóng gói nộp bài. \\
\bottomrule
\end{longtable}

\section{Thiết kế giao diện thử nghiệm}
Giao diện thử nghiệm được tổ chức thành ba tab. Tab “Phân tích Sentiment” là nơi nhập một review thật và dự đoán tiêu cực/trung lập/tích cực. Tab “Hỏi đáp RAG” là nơi đặt câu hỏi về review, xem hệ thống hiểu query, xem câu trả lời có bằng chứng và các review được truy xuất. Tab “Kỹ thuật \& đánh giá” hiển thị thông tin corpus, embedding model, FAISS index, baseline Precision@5 và đường dẫn artifact.

Thiết kế này phản ánh đúng đề tài: sentiment là phần phân tích cảm xúc chính, còn RAG giúp hỏi đáp và giải thích bằng review thật. Việc tách tab cũng tránh lỗi hiểu nhầm rằng sentiment model đang phân tích câu hỏi RAG.

\section{Backend và frontend}
Backend dùng FastAPI để load model sentiment, FAISS index, metadata, chạy retrieval và trả JSON cho frontend. Frontend được xây dựng bằng HTML, CSS và JavaScript, chịu trách nhiệm gọi API backend và render kết quả dưới dạng card, badge và panel kỹ thuật. Tên file cụ thể của phần hiện thực được liệt kê trong phụ lục.

\section{Nguyên tắc thiết kế}
Đề tài ưu tiên tính minh bạch. Với sentiment, báo cáo phải có metric và confusion matrix. Với RAG, câu trả lời phải có citation khớp với evidence cards. Với Module 6, kết quả benchmark phải ghi rõ là benchmark nhỏ, không thay thế đánh giá relevance thủ công lớn. Nguyên tắc này giúp báo cáo trung thực và tránh phóng đại năng lực hệ thống.

\section{Hợp đồng dữ liệu giữa các module}
Một điểm quan trọng trong thiết kế hệ thống là mỗi module cần có input/output rõ ràng. Module 1 xuất dữ liệu sentiment đã chuẩn hóa để Module 2 huấn luyện. Nếu tên cột text, rating, sentiment thay đổi không kiểm soát, Module 2 sẽ lỗi hoặc học sai. Tương tự, Module 4 cần corpus có trường retrieval text và metadata. Nếu metadata thiếu rating hoặc category, giao diện RAG vẫn có thể trả review nhưng không hiển thị đầy đủ evidence card.

Hợp đồng dữ liệu còn giúp debug. Khi giao diện không load được backend, có thể kiểm tra health endpoint. Khi RAG không trả kết quả, có thể kiểm tra trạng thái của FAISS index và metadata. Khi sentiment không chạy, có thể kiểm tra artifact model đã được load hay chưa. Cách chia module giúp lỗi được khoanh vùng nhanh hơn so với viết tất cả logic trong một notebook duy nhất.

\section{Thiết kế state management trong giao diện}
Trong phiên bản giao diện cuối, output RAG luôn ghi rõ “Kết quả cho câu hỏi nào”. Nếu người dùng sửa textarea sau khi đã có output cũ nhưng chưa bấm Chạy RAG, UI cần cảnh báo rằng câu hỏi đã thay đổi. Đây là chi tiết nhỏ nhưng quan trọng: nếu không có state management, người xem có thể nghĩ câu trả lời cũ là kết quả của câu hỏi mới.

Tab Sentiment và tab RAG cũng được tách riêng để tránh hiểu nhầm. Sentiment model phân tích review thật, không phân tích câu hỏi RAG. Trong RAG, hệ thống có thể suy luận intent/aspect từ query để áp dụng filter, nhưng đó là query understanding, không phải sentiment classification cho review.

\section{Thiết kế evidence card}
Evidence card cần hiển thị citation number, điểm liên quan, ngành hàng, số sao, sentiment và đoạn review ngắn. Nếu review quá dài, UI chỉ hiển thị excerpt để tránh phá layout. Citation trong answer phải khớp với số trên evidence card. Thiết kế này giúp người dùng kiểm tra câu trả lời nhanh: đọc summary trước, sau đó bấm/xem evidence nếu cần.

Các thông tin kỹ thuật như đường dẫn FAISS index, metadata path, candidates, returned, thời gian truy xuất được đưa vào vùng “Chi tiết kỹ thuật” hoặc tab kỹ thuật. Điều này giúp giao diện chính dễ đọc, nhưng vẫn đủ thông tin để kiểm tra pipeline.

\section{Các lựa chọn thiết kế thay thế}
Có thể xây dựng hệ thống theo hướng chỉ dùng notebook Colab. Cách đó nhanh cho thử nghiệm nhưng khó tách biệt các thành phần vận hành. Có thể dùng LLM API để sinh câu trả lời tự nhiên hơn, nhưng cần API key và có rủi ro hallucination. Có thể dùng PhoBERT cho sentiment ngay từ đầu, nhưng tăng độ phức tạp và thời gian train. Đề tài chọn kiến trúc hiện tại vì cân bằng giữa chất lượng, tính minh bạch, khả năng chạy local và khả năng tái lập kết quả.

